El segundo paso fue compilar y medir dentro de Docker. Un solo guion lo hace todo: baja, comprueba, construye la imagen, compila y mide. Tarda unas tres horas.

¿Por qué Docker? Porque nos da el mismo entorno que usaron los autores, Ubuntu 22.04 con g++-12, sin instalar nada en el sistema.

Docker no es una máquina virtual: no hay un procesador simulado. El programa corre en la misma CPU y con el mismo núcleo de Linux. Es como correr en la misma pista con un dorsal puesto.

Aun así, que Docker no cambie los números no se puede suponer. Lo comprobamos midiendo, y lo veremos en la parte 4.
